\documentclass[12pt]{article}
\usepackage[T1]{fontenc}
\usepackage[utf8]{inputenc}
\usepackage{lmodern}
\usepackage{textcomp}
\usepackage{enumitem}
\usepackage{geometry}
\geometry{margin=1in}
\begin{document}
\begin{center}
Answer Key - History - Graduate Level Assessment 17 \
\small (Answers are ordered exactly as in the question paper.)
\end{center}

\section*{Multiple Choice}
\begin{enumerate}[label=\arabic*.]
\item \textbf{Answer:} C
\\ \textit{Options:} A) monument for an Inca king.; B) hunting ground for Inca rulers.; C) relaxing getaway for elites.; D) military training base.; E) defensive stronghold.
\\ \textit{Note:} Machu Picchu was constructed as a luxurious retreat for Incan elites, serving both as a ceremonial site and a private residence, reflecting the social hierarchy and the importance of leisure among the ruling class in Inca society.
\item \textbf{Answer:} A
\\ \textit{Options:} A) a new type of metal alloy.; B) the earliest written documents.; C) the remains of a massive wooden ship believed to be Noah's Ark.; D) animal remains scattered on the surface of the ground.; E) the ruins of an ancient civilization under the sea.
\\ \textit{Note:} John Frere's discovery in 1797 of a new type of metal alloy was significant because it contributed to the understanding of prehistoric metallurgy and the development of new materials, which was a crucial advancement in both historical and archaeological contexts.
\item \textbf{Answer:} D
\\ \textit{Options:} A) dependence on megafauna.; B) hunting and gathering.; C) dependence on fur-bearing animals.; D) diversity.
\\ \textit{Note:} The correct answer is diversity because post-Pleistocene Europe experienced a variety of ecosystems and climates that supported a wide range of flora and fauna, leading to a rich and varied food resource base for human populations.
\item \textbf{Answer:} C
\\ \textit{Options:} A) Iberomaurusian culture; Maglemosian; B) Backed blade; microlith blade; C) Iberomaurusian culture; Capsian culture; D) Natufian; Iberomaurusian culture; E) Capsian culture; microlith blade
\\ \textit{Note:} The correct answer is supported by the fact that the Iberomaurusian culture, characterized by its unique stone tools and subsistence patterns, was gradually supplanted by the Capsian culture, which introduced new technological innovations and adaptations to changing environmental conditions in Holocene Africa.
\item \textbf{Answer:} A
\\ \textit{Options:} A) gas chromatography-mass spectrometry and atomic force microscopy; B) infrared spectroscopy and ultraviolet-visible spectroscopy; C) Raman spectroscopy and Fourier-transform infrared spectroscopy; D) electron microscopy and nuclear magnetic resonance; E) gravitational wave detection and atomic absorption spectroscopy
\\ \textit{Note:} Gas chromatography-mass spectrometry and atomic force microscopy are both advanced analytical techniques used to detect and quantify trace elements in various samples, making them suitable for trace element analysis.
\end{enumerate}

\section*{True / False}
\begin{enumerate}[label=\arabic*.]
\setcounter{enumi}{5}
\item \textbf{Answer:} True
\\ \textit{Options:} A) True; B) False
\\ \textit{Note:} According to the passage: Pergamum Altar
\item \textbf{Answer:} False
\\ \textit{Options:} A) True; B) False
\\ \textit{Note:} The correct answer is: 35\%
\item \textbf{Answer:} True
\\ \textit{Options:} A) True; B) False
\\ \textit{Note:} According to the passage: ESPN
\item \textbf{Answer:} False
\\ \textit{Options:} A) True; B) False
\\ \textit{Note:} The correct answer is: the cathedral of Serres
\item \textbf{Answer:} False
\\ \textit{Options:} A) True; B) False
\\ \textit{Note:} The correct answer is: aguas frescas and atole
\end{enumerate}

\section*{Long Form Response}
\begin{enumerate}[label=\arabic*.]
\setcounter{enumi}{10}
\item \textbf{Answer:} preserving tradition
\\ \textit{Note:} Expected answer: preserving tradition
\item \textbf{Answer:} maakond
\\ \textit{Note:} Expected answer: maakond
\end{enumerate}

\vfill
\noindent\textit{End of Answer Key}
\end{document}